% RH paper — § The saturated Selberg trace closure
% Drafted body introducing the formed-layer shell \Sel and its
% admissibility-as-hypothesis disclosure.
%
% statements-of-record rows resolved here: see
% paper/rh/notes/section-outline.md (this file's row) and
% paper/notes/statements-of-record.yml (target_section_hint field).

\section{The saturated Selberg trace closure}
\label{sec:sat_sel_shell}

With the functional-equation involution, the anti-invariant readout,
and the typed zero ledger in place, the formed-layer object of this
paper can be named.  It is the saturated completed Selberg trace
closure \(\Sel\).  In plain mathematical terms, \(\Sel\) is a typed
shell that packages, in one structure, the carriers, trace
instrument, quotient data, zero ledger, and audit data needed to apply
the duality-confinement mechanism to the Riemann-zeta substrate.  The
word ``saturated'' records the closure status of the shell; the phrase
``completed Selberg trace'' records the substrate being organized:
the completed zeta function together with the trace-relevant data
through which the later domination records are stated.

\begin{definition}[Saturated completed Selberg trace closure]
\label{def:rh:sat-sel-shell}
The saturated completed Selberg trace closure \(\Sel\) is a typed
shell consisting of the following thirteen fields:
\begin{enumerate}[label=\textup{(\arabic*)}]
  \item \(H_L\), the completed \(L\)-history carrier;
  \item \(\Itr\), a lawful trace instrument;
  \item \(\mathrm{EQ}_{\mathrm{tr}}\), the saturated trace-observable
        family;
  \item \(\mathrm{EM}_{\mathrm{zero}}\), the predictive verifier
        events for the zero-side data;
  \item \(Q_L\), the current quotient carrier \(H_L/{\sim_Q}\);
  \item \(M_L\), the predictive quotient carrier \(H_L/{\sim_M}\);
  \item \(\pi_L : M_L\to Q_L\), the comparison map between the two
        quotient sides;
  \item \(\JL\), the functional-equation involution from
        \Cref{def:rh:fe-involution};
  \item \(\Lambda_L\), the completed \(L\)-function, specialized here
        to \(\Lzeta\);
  \item \(\Znt\), the nontrivial-zero ledger from
        \Cref{def:rh:nontrivial-zero-ledger};
  \item \(\AZ\), the anti-invariant zero ledger from
        \Cref{def:rh:anti-invariant-zero-ledger};
  \item \(\mathrm{Vis}_{L}\), the instrument-relative visibility map;
  \item \(\mathrm{Audit}_{L}\), the audit provenance and admissibility
        field for the shell.
\end{enumerate}
The standing assumption attached to \(\Sel\) is that this shell is
admissible for the relevant Foundations-II formation schemas.
\end{definition}

The field list separates the three kinds of data that will be used
below.  The first cluster, \(H_L,\mathrm{EQ}_{\mathrm{tr}},
\mathrm{EM}_{\mathrm{zero}},Q_L,M_L\), and \(\pi_L\), records the
history and quotient architecture of the formed closure.  These fields
say which typed carrier is being studied, which trace observables are
available, which verifier events have been retained, and how the
current and predictive quotient views are compared.  They are not a
new analytic construction of the zeta function; they are the
formed-layer bookkeeping that lets the trace mechanism refer to the
same substrate throughout the paper.

The second cluster contains the load-bearing RH objects from
\Cref{sec:involution_and_ledger}.  The involution \(\JL\) supplies the
critical-line symmetry.  The completed function \(\Lambda_L=\Lzeta\)
names the zeta substrate whose nontrivial zeros populate \(\Znt\).
The ledger \(\Znt\) carries the multiplicity-weighted zero data, and
\(\AZ\) is the nonnegative scalar built from the signed readout
\(\psimRH\).  These are the fields through which the later translation
theorem speaks: \(\AZ=0\) is the trace-side statement that will be
converted into the critical-line statement for all zeros in \(\Znt\).

The third cluster, \(\Itr,\mathrm{Vis}_{L}\), and
\(\mathrm{Audit}_{L}\), is the instrument and audit layer.  The trace
instrument \(\Itr\) is the lawful trace functional attached to the
positive-cone side of the shell.  The visibility map records which
instrument-relative observations are available at the formed layer.
The audit field records the provenance of the construction-time data
and carries the admissibility status of the shell.  This is the
precise sense in which \(\Sel\) is the object on which the conditional
landing chain operates: it is not merely a list of zeros, but a typed
trace shell in which the zero ledger, trace instrument, and audit
status coexist.

An important encoding disclosure belongs to this definition.  The
admissibility of \(\Sel\) with respect to the seven Foundations-II
schemas is encoded as a Foundations-II-level hypothesis carried by
\(\mathrm{Audit}_{L}\) \cite{TsiokosFoundationsII2026}.  It is not
derived in this paper from a separate construction of the completed
Selberg trace instance.  A substantive formal derivation of
admissibility would require unfolding those seven schemas as explicit
predicates and verifying each one for the completed-zeta trace shell.
For the present paper, admissibility is a standing assumption of the
typed shell, recorded at construction time in the audit field.

The Selberg trace formula context behind the terminology is the
classical trace-formula machinery associated with Selberg's original
work and later expositions such as Iwaniec's \emph{Spectral Methods
of Automorphic Forms} and the Iwaniec--Kowalski analytic number
theory text.  Those classical references belong to the external
bibliography curation stage.  Here the role of that context is more
specific: the paper uses the typed shell above to record exactly what
the duality-confinement mechanism needs from the Selberg trace setting
without re-developing the full analytic trace-formula apparatus.

The next section proves translation theorem T: \(\AZ=0\) is
equivalent to the RH statement on the typed zero ledger of \(\Sel\).
With that translation in place and the recognition source supplied in
\Cref{sec:recognition_source}, the conditional landing chain can apply
the duality-confinement master theorem from the companion paper.
